\section{Introduction}
A target moving toward a camera grows in the image during exposure. Different image points blur by different amounts, so a single translation-invariant convolution generally cannot describe the blur. A hand push also varies in speed. Measuring depth during exposure could give a better estimate of this changing magnification, while coded illumination changes the contribution of each target position to the photograph.

We ask whether measured depth improves reconstruction of the same coded image compared with a straight line between the same exposure-endpoint depths. The comparison holds the image, illumination history, endpoint geometry, solver and scoring region fixed. We also compare each reconstruction with the blurred input to check whether a better motion estimate actually improves the image.

We built a manual guide path and used an existing upright target stand. The camera, depth board and mounting plate came from the lab kit. Monitor illumination and concurrent image, depth and command logs let us model motion within a conventional long exposure. Our processing aligns these records, constructs a scale-blur operator and estimates the sharp image. All motion was produced by hand. This report presents the supplementary physical dataset and every primary outcome; the earlier five-trial pilot used a different protocol and is kept separate.
